\documentclass[10pt,a4paper]{amsart}
\usepackage[margin=19mm]{geometry}
\usepackage{amsmath,amssymb,mathtools}
\usepackage[scr=rsfs,cal=euler]{mathalfa}
\usepackage{xfrac}
\usepackage[unicode,hidelinks,bookmarks=false]{hyperref}
\newcommand{\ie}{i.e.\ }
\newcommand{\cf}{cf.\ }
\newcommand{\resp}{respectively\ }
\newcommand{\Subsection}{\subsection}
\providecommand{\nobreakdash}{\nobreak\hbox{-}}
\setlength{\emergencystretch}{2em}
\multlinegap0pt
\usepackage{bm}
\usepackage{xcolor}
\usepackage{mathtools}
\usepackage{amssymb}
\usepackage{enumerate}
\usepackage[shortlabels]{enumitem}
\newtheorem{theorem}{Theorem}
\newtheorem{proposition}{Proposition}
\newcommand{\R}{\mathbb{R}}
\title{Time-optimal neural feedback control of nilpotent systems as a binary classification problem}
\author{Sara Bicego, Samuel Gue, Dante Kalise, Nelly Villamizar}
\date{Source: arXiv:2503.17581v2 (CC BY 4.0)}
\begin{document}
\maketitle

\noindent\emph{Source and licence.} Time-optimal neural feedback control of nilpotent systems as a binary classification problem. Version arXiv:2503.17581v2 (CC BY 4.0), distributed by arXiv under the Creative Commons Attribution 4.0 International licence, http://creativecommons.org/licenses/by/4.0/, as recorded in the arXiv licence metadata for this exact version and shown on the abstract page https://arxiv.org/abs/2503.17581v2. Full licence notice: \texttt{licenses/arxiv-2503.17581-CC-BY-4.0.txt}.\par\smallskip

\noindent\textit{Editorial scope.} Counted unit: the article's proposition prop:21 (source0.tex lines 106--113). The article's complete original proof of it is delivered with it (source0.tex lines 114--160, one \texttt{proof} environment of 47 lines). No mathematics is written, repaired or re-derived: the statement and the proof are the article's own lines, in the article's own notation, and no retained line is substituted or re-flowed.  Retained uncounted as the source-local context the proof needs: lines 83--90; lines 93--94.  Nothing was omitted from any retained span, so the package carries no omission record.  The retained lines cite 1 works, whose entries are transcribed verbatim from the article's own bibliography source in the e-print and delivered with short editorial labels recorded in the item's reference mapping.  No retained line contains a figure environment, an image-inclusion command or drawing code, so no image is delivered and the package declares no asset dependency.  Editorial formatting only, in addition: an AMS \texttt{amsart} preamble that restates the article's own declarations and macros used by the retained lines, namely the environments \texttt{theorem}, \texttt{proposition} on separate counters, together with the standard packages those lines need.  The retained environments print in this document as Theorem 1, Proposition 1 in order of appearance, whereas the article prints the same items with the numbering of its own full document.  The complete native source of the article as distributed in the arXiv e-print for this exact version is bundled unchanged as \texttt{sources/175193/source0.tex}.
\begin{align}\label{eq:linearSystem}
 \underset{u\in\mathcal{U}}{\min} \;\;T\quad 
 \text{subject to} \quad
 &\begin{cases}\dot{\bm{y}}(t)=A\bm{y}(t)+\bm{b}u(t)\\
 \bm{y}(0)=\bm{x}\\
 \bm{y}(T)=\bm{0},
 \end{cases}
\end{align}

\begin{theorem}[Chapter III, \cite{PMP}]\label{switching} Assume that Kalman's controllability condition holds for the pair (A,$\bm{b}$), and that all the eigenvalues of $A$ are real. Then, for every initial condition $\bm{x}$, there exists a unique time-optimal control signal $u$. This control signal is bang-bang, i.e. taking values in $\{-1,1\}$, with at most $n-1$ switches. 
\end{theorem}

\begin{proposition}\label{prop:21}
 Let $A$ be an $n\times n$ nilpotent Jordan block and $\bm{b}\in\R^n$ satisfying the assumptions in Theorem \ref{switching}. Given an initial condition $\bm{x}$, the non-negative increments $\{t_i\}_{i=1}^n$ such that $T=t_1+\cdots+t_n$, associated to the time-optimal control signal $u(t)=u_i$, for all $t \in \Omega_i$, with $u_i=\pm 1$, satisfy the polynomial system
 \begin{equation}\label{eq:polynomialSystem2}
 0=-u(0){\sum_{k=1}^{n-i+1} b_{k+i-1}}\sum_{\substack{\alpha_1,\ldots,\alpha_n\geq0\\
 \alpha_1+\cdots+\alpha_n=k}}(-1)^{\max \{j\colon \alpha_j\neq0\}}\frac{t_1^{\alpha_1}\cdots t_n^{\alpha_n}}{\alpha_1!\cdots\alpha_n!}+{\sum_{k=i}^{n}\frac{T^{k-i}}{(k-i)!}x_{k}},
 \end{equation}
 for $i=1,\dots,n,$ where $b_j, x_j$ denote the coordinates of $\bm b$ and $\bm x$, respectively.
 \end{proposition} 

 \begin{proof}
 First, recall that for a given control signal, the trajectory associated to the linear system in \eqref{eq:linearSystem} is given by 
 \begin{equation*}
 \bm{y}(t)=e^{tA}\bm{y}(0)+e^{tA}\int_{0}^{t}e^{-\tau A}\bm{b}\,u(\tau)\,d\tau. 
 \end{equation*} 
Evaluating at $T=t_1+\cdots+t_n$, and imposing that $\bm{y}(T)=\bm{0}$ implies
 \[\bm{0}=e^{TA}\bm{x}+e^{TA}\int_{0}^{T}e^{-\tau A}\bm{b}\,u(\tau)\, d\tau,
 \]
 and integration by parts yields
 \begin{align}
 \bm{0}&=e^{AT}\bm{x}+e^{TA}\sum_{k=1}^{n} \bigl(A^{k-1}\, e^{-T A}\bm b\, U_k(T)-A^{k-1}\bm b\,U_k(0)\bigr)=e^{TA}\bm{x}+\sum_{k=1}^{n} A^{k-1}\,\bm{b}\,U_k(T),\label{eq:eq}
 \end{align}
 where 
 \[U_{k}(t)=\int_0^t U_{k-1}(\tau)\, d\tau, \text{ for\; } i\geq 1, \text{ with\; } U_1(t)=\int_{0}^t u(\tau) d\tau.
 \]
 Since $u$ is piecewise constant on the partition of $\Omega$ determined by the $t_i$'s, 
 \[U_1(T)=u(0)\sum_{i=1}^{n}\int_{T_{i-1}}^{T_i} (-1)^{i-1}d\tau=u(0)(t_1-t_2+\dots(-1)^{n-1}t_n),\] 
 where, as above, $T_i=t_1+\cdots+t_i$. 
 For any $1\leq k\leq n-1$, we have
 \begin{align*} 
 U_{k}(T)&=\int_{0}^{T}U_{k-1}(\tau)d\tau=\sum_{i=1}^{n}\int_{T_{i-1}}^{T_i}U_{k-1}(\tau)d\tau
 =
 u(0)\sum_{i=1}^n\int_{T_{i-1}}^{T_i} (-1)^{i-1}\frac{\tau^{k-1}}{(k-1)!}d\tau
 \\
 &= u(0)\sum_{i=1}^n (-1)^{i-1}\biggl[\frac{(t_i+T_{i-1})^k}{k!}-\frac{(T_{i-1})^k}{k!}\biggr]= u(0)\sum_{i=1}^n\sum_{j=1}^k\frac{(-1)^{i-1}}{(k-j)!j!}t_i^jT_{i-1}^{k-j}.\nonumber
 \end{align*}
 Since
 \[\left(t_1+\cdots+t_{i-1}\right)^{k-j}
 =
 \sum_{\substack{\alpha_1+\cdots+\alpha_{i-1}=k-j\\ \alpha_{1},\dots,\alpha_{i-1}\geq 0}}\frac{(k-j)!}{\alpha_{1}! \alpha_{2}!\cdots \alpha_{i-1}!}t_1^{\alpha_1} \cdots t_{i-1}^{\alpha_{i-1}},
 \]
 then
 \begin{align}
 U_{k}(T)&=
 u(0)\sum_{i=1}^n\sum_{\substack{\alpha_1+\cdots+\alpha_{i}=k\\\alpha_i,\dots,\alpha_{i-1}\geq 0,\,\alpha_i>0}}\frac{(-1)^{i-1}}{\alpha_{1}!\cdots \alpha_{i}!}t_1^{\alpha_1} \cdots t_{i}^{\alpha_{i}}\nonumber\\
 &=
 -u(0)\sum_{\substack{\alpha_1+\dots+\alpha_n=k\\\alpha_1,\ldots,\alpha_n\geq0}}\frac{(-1)^{\max\{j\colon \alpha_j\neq 0\}}}{\alpha_{1}! \cdots \alpha_{n}!}t_1^{\alpha_1}\cdots t_n^{\alpha_n}.\label{eq:sum}
 \end{align}
 On the other hand, since $A$ is a nilpotent Jordan block, we have $A_{ij}=\delta_{i,j-1}$, for $i,j=1, \dots, n$.
 Then, for $1\leq k\leq n$ we have $(A^{k-1})_{i,j}=\delta_{i,i+k-1}$. Thus, $\bigl( A^{k-1}\bm{b}\bigr)_{i}=b_{k+i-1}$ for $i=1,\dots, n-k+1$, and $\bigl( A^{k-1}\bm{b}\bigr)_{i}=0$ otherwise.
 
 Finally, replacing \eqref{eq:sum} in \eqref{eq:eq}, and noting that the $i$-th component of $e^{TA}\bm{x}$ is given by
 \begin{equation*}
 (e^{TA}\bm{x})_i=\sum_{k=1}^{n}(e^{TA})_{i,k}{x}_k=\sum_{k=i}^{n}\frac{T^{k-i}}{(k-i)!}{x}_k,
 \end{equation*} 
 we obtain the equality \eqref{eq:polynomialSystem2}, as required.
 \end{proof}

\begin{thebibliography}{99}
\bibitem[PMP]{PMP} L.~Pontryagin, V.~Boltyanskii, R.~Gamkrelidze, and E.~Mishchenko. \newblock {\em The mathematical theory of optimal processes}. \newblock A Pergamon Press Book. The Macmillan Company, New York, 1964. \newblock Translated by D. E. Brown.
\end{thebibliography}
\end{document}
